\begin{table}[H]
\centering
\caption{\textbf{Genomic geometry improves prediction of the edited-site fraction in all four tissues.} Larger inter-arm gaps correlate with lower burden (Spearman $\rho$). Burden is the high-editing fraction among eligible high/low sites within a pair. Adding geometry to graph-derived summaries improves pair-level $R^2$; $\Delta R^2$ is relative to graph summaries alone, with a 95\% pair-clustered bootstrap interval. ``Any high'' asks whether a pair contains at least one site with EditingIndex $\geq15\%$. These outcomes measure editing burden, not molecular realization or ADAR engagement.}\label{tab:substraterealization}
\begin{tabular}{lrrrrr}
\toprule
Context & Gap--burden $\rho$ & Graph-only $R^2$ & Combined $R^2$ & $\Delta R^2$ [95\% CI], points & Any-high AUROC \\
\midrule
Artery & -0.482 & 0.053 & 0.270 & +21.64 [+16.39, +26.93] & 0.816 \\
Brain & -0.425 & 0.050 & 0.239 & +19.11 [+13.94, +23.68] & 0.749 \\
Liver & -0.475 & 0.070 & 0.277 & +20.82 [+15.75, +26.03] & 0.728 \\
Skeletal muscle & -0.334 & 0.010 & 0.050 & +4.12 [+0.99, +6.98] & 0.705 \\
\bottomrule
\end{tabular}
\end{table}
